\documentclass{article}
% --- ENCODING & LANGUAGE (MUST BE FIRST) ---
\usepackage[T1]{fontenc}
\usepackage[utf8]{inputenc}
\usepackage[swedish]{babel}
\usepackage[margin=2cm]{geometry}

% --- CORE MATH & SYMBOLS ---
\usepackage{amsmath, amssymb, amsthm}

% --- GRAPHICS & UTILS ---
\usepackage{graphicx}
\usepackage{xcolor}
\usepackage{ulem}
\usepackage{changepage}
\usepackage{listings}
\lstset{language=Python}

% --- NATURE COLOR DEFINITIONS ---f
\definecolor{natyellow}{HTML}{D4AC0D}   % Teorem, Lemma, Proposition, Corollary
\definecolor{natblue}{HTML}{1A5276}    % Definition
\definecolor{natorange}{HTML}{CA6F1E}   % Proof
\definecolor{natgray}{HTML}{696969}    % Remark
\definecolor{natbrown}{HTML}{6E2C00}    % Ex
\definecolor{natgreen}{HTML}{27AE60}    % Keyword
\definecolor{pink}{HTML}{d951bb} 

% --- SÄKERT HJÄLPKOMMANDO FÖR RUBRIKER ---
\newcommand{\naturetitle}[3]{%
  \par\addvspace{0.8em}%
  \noindent%
  \colorbox{#1!25}{\textbf{\textcolor{black}{#2}}}\textbf{ (#3):}%
  \vspace{0.3em}\noindent\ignorespaces
}

% --- MILJÖER ---

% 1. Definition environment
\newenvironment{defi}[1]%
{\naturetitle{natblue}{Definition}{#1}\par}%
{\addvspace{0.8em}}

% 2. Theorem, Lemma, Proposition, Corollary (Gula)
\newenvironment{theorem}[1]%
{\naturetitle{natyellow}{Theorem}{#1}\par}%
{\vspace{0.8em}}

\newenvironment{lemma}[1]%
{\naturetitle{natyellow}{Lemma}{#1}\par}%
{\vspace{0.8em}}

\newenvironment{proposition}[1]%
{\naturetitle{natyellow}{Proposition}{#1}\par}%
{\vspace{0.8em}}

\newenvironment{corollary}[1]%
{\naturetitle{natyellow}{Corollary}{#1}\par}%
{\vspace{0.8em}}

% 3. Proof (Orange)
\newenvironment{prf}[1]%
{\naturetitle{natorange}{Proof}{#1}\par\itshape}%
{\hfill$\square$\vspace{0.5em}}

% 4. Ex / Example (Brun)
\newenvironment{example}[1]%
{\naturetitle{natbrown}{Example}{#1}\par}%
{\vspace{0.5em}}

% 5. Remark (Grå - med högerpil och indrag)
\newenvironment{remark}
{\par\addvspace{0.05em}\noindent\textbf{\textcolor{natgray}{$\hookrightarrow$ \uline{Remark:}}}\\[0.1em]\small\itshape\begin{adjustwidth}{0.1em}{0.1em}}
{\end{adjustwidth}\par\addvspace{0.05em}}

% 6. Keyword (Grön)
\newcommand{\keyword}[1]{\textbf{\textcolor{natgreen}{#1}}}

% 7. Method 
\newenvironment{method}[1]%
{\naturetitle{pink}{Method}{#1}\par}%
{\addvspace{0.8em}}

\title{Rapport}
\author{Disa Degerholm}
\date{\today}

\begin{document}

\maketitle
\thispagestyle{empty}

\section*{Uppgift 3. Existens men inte entydighet}

\subsection*{Uppgift 3. Disa Degerholm}

Betrakta begynnelsevärdesproblemet
\begin{equation}
    y' = \sqrt{|y|}, \quad y(0) = 0.
\end{equation}

\subsection*{1. Analytisk lösning och icke-entydighet}
Picards existens- och entydighetssats (Simmons kap.~13) kräver att $f(y) = \sqrt{|y|}$ är lokalt Lipschitzkontinuerlig i en omgivning av $y = 0$, dvs.~att differenskvoten är begränsad:
\begin{equation}
    \lim_{y \to 0^+} \frac{|f(y) - f(0)|}{|y - 0|} = \lim_{y \to 0^+} \frac{\sqrt{y}}{y} = \lim_{y \to 0^+} \frac{1}{\sqrt{y}} = \infty.
\end{equation}
Eftersom Lipschitzvillkoret fallerar i $y = 0$ upphör entydigheten (Peanos existenssats garanterar endast existens).

Separation av variabler för $y \ne 0$:
\begin{itemize}
    \item För $y > 0$: $\int y^{-1/2}\,dy = \int dt \implies 2\sqrt{y} = t - c_2 \implies y(t) = \frac{1}{4}(t - c_2)^2$ ($t \ge c_2$).
    \item För $y < 0$: $\int (-y)^{-1/2}\,dy = \int dt \implies -2\sqrt{-y} = t - c_1 \implies y(t) = -\frac{1}{4}(t - c_1)^2$ ($t \le c_1$).
\end{itemize}
Eftersom $y' \ge 0$ är $y(t)$ icke-avtagande och kan ligga kvar på jämvikten $y \equiv 0$ under ett godtyckligt tidsintervall $[c_1, c_2]$:
\begin{equation}
    y(t) = \begin{cases}
        -\frac{1}{4}(t - c_1)^2, & t \le c_1, \\
        0, & c_1 < t < c_2, \\
        \frac{1}{4}(t - c_2)^2, & t \ge c_2,
    \end{cases}
    \quad \text{där } c_1 \le 0 \le c_2.
\end{equation}
Detta ger en tvåparametrig familj av oändligt många lösningar.

\subsection*{2. Numerisk lösning med RK4}
Med $y_0 = 0$ fås i första RK4-steget:
\begin{equation}
    k_1 = f(0) = 0 \implies k_2 = k_3 = k_4 = 0 \implies y_1 = y_0 = 0.
\end{equation}
Genom induktion gäller $y_n = 0$ för alla $n$. En standardlösare förblir deterministiskt låst i den stationära lösningen $y(t) \equiv 0$ och kan inte numeriskt detektera förgreningen utan yttre störning.

\subsection*{3. Hur länge lösningen fastnar vid $y(-1) = -1$}
För $y < 0$ ger separation av variabler med $y(-1) = -1$:
\begin{equation}
    -2\sqrt{-y} = t + C \implies -2\sqrt{1} = -1 + C \implies C = -1,
\end{equation}
vilket ger $y(t) = -\frac{1}{4}(1 - t)^2$ för $t \le 1$. Analytiskt nås $y = 0$ vid $t = 1$, varpå lösningen kan kvarstanna i $y \equiv 0$ godtyckligt länge.

\begin{itemize}
    \item \textbf{Tid till tröskeln $|y| \le 10^{-6}$:}
    För $y < 0$ ger separation av variabler med $y(-1) = -1$:
    \[
    \int (-y)^{-1/2}\,dy = \int dt \implies -2\sqrt{-y} = t - 1 \implies y(t) = -\frac{1}{4}(1 - t)^2.
    \]
    Tröskelvillkoret $|y| \le 10^{-6}$ ger:
    \[
    \frac{1}{4}(1 - t)^2 \le 10^{-6} \implies |1 - t| \le 2 \cdot 10^{-3} \implies t \approx 0{,}998.
    \]
    Tiden från startpunkten $t = -1$ blir därmed $\Delta t = 0{,}998 - (-1) = 1{,}998$.

    \item \textbf{Numeriskt beteende:}
    \begin{itemize}
        \item Om $y_n = 0$ exakt blir alla lutningar noll:
        \[
        k_1 = f(t_n, 0) = \sqrt{|0|} = 0 \implies k_2 = k_3 = k_4 = 0 \implies y_{n+1} = y_n = 0,
        \]
        vilket låser lösaren permanent i jämvikten.
        \item Om avrundningsfel ger $y_n = \delta > 0$ fås efter ett Euler-steg (eller motsvarande för RK4):
        \[
        y_{n+1} = \delta + h\sqrt{\delta} \implies \frac{y_{n+1}}{y_n} = 1 + \frac{h}{\sqrt{\delta}}.
        \]
        Eftersom $\lim_{\delta \to 0^+} \frac{h}{\sqrt{\delta}} = \infty$ växer tillskottet $\Delta y = h\sqrt{\delta} \gg \delta$ explosionsartat, varför lösningen lämnar intervallet $|y| \le 10^{-6}$ på ett fåtal steg.
    \end{itemize}
\end{itemize}

\subsection*{Uppgift 3. Hampus Hillding}

\subsection*{Uppgift 3. Skillnader och likheter}

\begin{enumerate}
    \item \textbf{Skillnader för härledning till ODE:n:} Disa tittade på Lipschitz kontinuitet i punkten y=0 medan Hampus löste ODE:n direkt. Detta gjorde att när Hampus tittade på sin numeriska lösning antog funktionen 0 inte bara i ett värde utan flera medan Disa hittade det.
    \item \textbf{Lös problemet numeriskt med kommentarer:} Både Disa och Hampus fick att $y=0, \forall x\in \mathbb{R}$ oavsett om man använde RK4 eller Euler. Det var annorlunda än Hampus analytiska lösning.
    \item \textbf{Hur länge fastnar den numeriska lösningen:} Disa trodde att om $y_n$ träffas exakt på 0 så kommer derivatan bli 0 och lösaren fastnar i $y=0$, hon trodde detta möjligen kunde ske om metoder med lägre ordning skulle användas tex Eulers metod. Men då visade Hampus att den faktiskt tillslut ändå växte förbi 0. Enligt Disas analytiska lösning så blev tiden ca 1 tidsenhet på $y=0$ enligt hennes härledning. ??? (när nås bara y=0)
\end{enumerate}

\section*{Uppgift 4. Titel}

\subsection*{Uppgift 4. Disa Degerholm}

Betrakta begynnelsevärdesproblemet
\begin{equation}
    y' = y^2, \quad y(0) = y_0 = \frac{N}{100}.
\end{equation}

\subsection*{1. Analytisk lösning och singularitet}
För $y \neq 0$ ger separation av variabler med $y(0) = y_0$:
\begin{equation}
    \int y^{-2}\,dy = \int dt \implies -\frac{1}{y} = t - \frac{1}{y_0} \implies y(t) = \frac{y_0}{1 - y_0 t} = \frac{N}{100 - N t}.
\end{equation}
Lösningen har en vertikal asymptot (smäller) då nämnaren försvinner:
\begin{equation}
    t^* = \frac{1}{y_0} = \frac{100}{N}.
\end{equation}
För $N = 82$ är den exakta singularitetstiden $t^* = \frac{100}{82} \approx 1{,}219512$.

\subsection*{2. Numerisk singularitet och skattning av $t^*$}
\begin{itemize}
    \item \textbf{Numeriskt bevis för singularitet:} 
    Vid numerisk integration (t.ex.\ RK4) växer $y_n$ och derivatan $y'_n = y_n^2$ explosionsartat då $t \to t^*$. Eftersom lösningskurvan blir vertikal kan ett fixt steg $h$ inte följa funktionen. Värdena överskrider snabbt gränsen för standardflyttal (\texttt{float64} $\sim 1{,}8 \cdot 10^{308}$) och genererar \texttt{inf}/\texttt{NaN}, vilket numeriskt påvisar singulariteten.

    \item \textbf{Uppskattning av $t^*$:}
    Genom att integrera framåt med steglängd $h = 0{,}01$ fram till ett tröskelvärde (t.ex.\ $y > 10^5$) avbryts beräkningen vid sista steget:
    \[
    t_{\text{stopp}} \approx 1{,}2100.
    \]
    Detta ger en direkt numerisk undre uppskattning av singularitetstiden som stämmer väl överens med det analytiska värdet $t^* = \frac{100}{82} \approx 1{,}2195$.
\end{itemize}

\subsection*{3. Bevis: $y(0) < 0 \implies y(t) < 0$ för alla $t > 0$}
Högerledet $f(y) = y^2$ har kontinuerlig partiell derivata $\frac{\partial f}{\partial y} = 2y$, vilket gör $f$ lokalt Lipschitzkontinuerlig. Enligt Picards entydighetssats (Simmons kap.~13) kan två skilda lösningskurvor aldrig skära varandra:
\begin{itemize}
    \item Funktionen $y(t) \equiv 0$ är en stationär lösning ($0' = 0^2$).
    \item Om det funnes ett $t_1 > 0$ med $y(t_1) \ge 0$, måste $y(t)$ enligt satsen om mellanliggande värden anta $y(t_0) = 0$ för något $t_0 \in (0, t_1]$.
    \item I punkten $(t_0, 0)$ skulle banan skära $y \equiv 0$, vilket motsäger Picards entydighet. Alltså gäller $y(t) < 0$ för alla $t > 0$.
\end{itemize}

\subsection*{4. Numeriskt felsteg och automatlösare}
Betrakta $y' = C y^2$ med $y_0 = \epsilon < 0$. Ett Euler-steg med steglängd $h$ ger:
\begin{equation}
    y_1 = \epsilon + h C \epsilon^2 = \epsilon (1 + h C \epsilon).
\end{equation}
Eftersom $\epsilon < 0$ blir $y_1 > 0$ om faktorn växlar tecken:
\begin{equation}
    1 + h C \epsilon < 0 \implies h > \frac{1}{C |\epsilon|}.
\end{equation}
För t.ex.\ $C = 50$, $\epsilon = -0{,}10$ och $h = 0{,}5$ ger RK4 ett falskt omslag till $y_1 = 0{,}0704 > 0$, varpå lösningen exploderar mot $+\infty$ istället för att avta mot noll.

\paragraph{Automatlösarens skydd:}
\begin{itemize}
    \item \textbf{Adaptiv steglängd:} Felet uppskattas via inbäddade metoder (t.ex.\ Runge--Kutta--Fehlberg) och $h$ skalas ned så att $h \ll \frac{1}{C|y|}$.
    \item \textbf{Invarians- och teckenkontroll:} Steget förkastas omedelbart om $y_n$ byter tecken ($y_n \cdot y_{n+1} \le 0$) eftersom teckenbyte är analytiskt förbjudet.
\end{itemize}

\subsection*{Uppgift 4. Hampus Hillding}

\subsection*{Uppgift 4. Skillnader och likheter}

\begin{enumerate}
    \item \textbf{Lös analytiskt mha personnummer:} Disa löste ODE:n för generellt N och stoppade sedan in sina sista två personnummersiffror. Hampus löste med sina direkt i formeln men båda fick samma svar (om N är samma).
    \item \textbf{När smäller det numeriskt?} Hampus använde Euler igen och Disa använde RK4. Detta gjorde att Disa fick fler iterationer innan overflow (??? why). Det "smällde" ungefär vid 1.1 för både Disa och Hampus vilket stämde överens med det analytiska. Vi ser att det "smäller" då programmet går in i "overflow".
    \item \textbf{Visa att om $y(0)<0$ så måste $y(t)<0$ för alla $t>0$:} Både DIsa och Hampus hade samma argument med att f är lokalt Lipschitz kontineruligt överallt och att då y(t)=0 är en lösning kan ingen annan lösning korsa den och blir därmed bara över eller under 0 i sitt y värde.
    \item \textbf{Numeriska exempel:} Hampus testade sig fram till ett stort nog C då han tänkte att ett stort C skulle ge ett $y'$ som var tillräckligt stor hopp för att komma över gränsen. Disa härledde formeln för det istället och kom fram till att $h > \frac{1}{C |\epsilon|}$.
    \item \textbf{Automatlösare:} Disa tänkte att en adaptiv steglängdskontroll (???) och detektering av teckenbyte (???) skulle kunna hjälpa. Hampus hade samma ideer.
\end{enumerate}

\section*{Uppgift 5. Titel}

\subsection*{Uppgift 5. Disa Degerholm}

Betrakta begynnelsevärdesproblemet
\begin{equation}
    y' = y, \quad y(0) = N.
\end{equation}

\subsection*{1. Analytisk lösning och globalt fel för RK4}
Differentialekvationen har den exakta lösningen:
\begin{equation}
    y(t) = N e^t \implies y(2^k) = N e^{2^k}.
\end{equation}
För $y' = y$ gäller $y^{(v)}(\xi) = N e^\xi$. Det lokala trunkeringsfelet per steg ges av:
\begin{equation}
    \epsilon_{\text{lokal}} = -\frac{y^{(v)}(\xi)}{180} h^5 = -\frac{N e^\xi}{180} h^5.
\end{equation}
Över $M = \frac{T}{h}$ steg fram till $T = 2^k$ begränsas det globala felet av:
\begin{equation}
    |E_M| \le M \cdot \max_{0 \le \xi \le T} |\epsilon_{\text{lokal}}(\xi)| \cdot e^{T - \xi} \le \frac{T}{h} \left(\frac{N e^T}{180} h^5\right) = \frac{N \cdot 2^k e^{2^k}}{180} h^4.
\end{equation}
Kravet på två korrekta decimaler ($|E_M| \le 0{,}005$) kräver:
\begin{equation}
    \frac{N \cdot 2^k e^{2^k}}{180} h^4 \le 0{,}005 \implies h \le \left(\frac{0{,}9}{N \cdot 2^k}\right)^{1/4} e^{-2^{k-2}}.
\end{equation}
Steglängden $h$ måste alltså minskas exponentiellt, vilket gör att antalet steg $M = \frac{2^k}{h} \propto 2^{5k/4} e^{2^{k-2}}$ växer explosionsartat.

\subsection*{2. Datorns begränsningar (Varför datorn ger upp)}
I standard flyttalsformat (\texttt{float64}, $\epsilon_{\text{mach}} \approx 2{,}22 \cdot 10^{-16}$) stoppas beräkningen av två skilda orsaker:
\begin{enumerate}
    \item \textbf{Förlust av precision ($k \ge 5$):} 
    Avståndet mellan två på varandra följande representerbara flyttal vid storleken $y$ är $\text{ulp} \approx \epsilon_{\text{mach}} \cdot y$. För två korrekta decimaler krävs $\epsilon_{\text{mach}} \cdot y \le 0{,}005$:
    \begin{equation}
        y \le \frac{0{,}005}{2{,}22 \cdot 10^{-16}} \approx 2{,}25 \cdot 10^{13} \implies 2^k \le \ln\left(\frac{2{,}25 \cdot 10^{13}}{N}\right).
    \end{equation}
    För $N = 82$ passeras gränsen redan vid $t \approx 26{,}3$, dvs.\ för alla $k \ge 5$ ($t_5 = 32$). Då är avståndet mellan talen större än toleransen, vilket gör två korrekta decimaler  omöjliga oavsett hur litet $h$ väljs då $|E_M| \le 0{,}005$ inte kan uppnås.

    \item \textbf{Aritmetiskt överflöde / \texttt{overflow} ($k \ge 10$):} 
    Det maximala talet i \texttt{float64} är $y_{\max} \approx 1{,}79 \cdot 10^{308}$. Detta överskrids då:
    \begin{equation}
        N e^{2^k} > 1{,}79 \cdot 10^{308} \implies 2^k > \ln\left(\frac{1{,}79 \cdot 10^{308}}{82}\right) \approx 705{,}4.
    \end{equation}
    Eftersom $t_9 = 512$ och $t_{10} = 1024$, slår beräkningen i hårdvarutaket vid $k = 10$ och returnerar \texttt{inf}.
\end{enumerate}

\subsection*{3. Resultat}
\begin{center}
\begin{tabular}{c r r c l}
\hline
$k$ & $t_k = 2^k$ & $y_{\text{exakt}}$ & Absolut fel & Status \\
\hline
$1$ & $2$ & $605{,}9026$ & $1{,}59 \cdot 10^{-6}$ & OK ($100$ steg) \\
$2$ & $4$ & $4477{,}0483$ & $3{,}70 \cdot 10^{-4}$ & OK ($100$ steg) \\
$3$ & $8$ & $244438{,}5549$ & $2{,}56 \cdot 10^{-3}$ & OK ($400$ steg) \\
$4$ & $16$ & $7{,}29 \cdot 10^8$ & $3{,}79 \cdot 10^{-3}$ & OK ($6400$ steg) \\
$5$ & $32$ & $6{,}47 \cdot 10^{15}$ & -- & Stopp: $\text{ulp} > 0{,}005$ \\
$10$ & $1024$ & $\infty$ & -- & Krasch: \texttt{overflow} (\texttt{inf}) \\
\hline
\end{tabular}
\end{center}

\subsection*{Uppgift 5. Hampus Hillding}

\subsection*{Uppgift 5. Skillnader och likheter}

\begin{enumerate}
    \item \textbf{Lös numeriskt med 2 korrekta decimaler:} Hampus gjorde Euler, Disa gjorde RK4. Både tog fram den teoretiska formeln för det totala trunkeringsfelet $|E_n|$. 
    \item \textbf{Jämför teoretiska och numeriska felen:} ??? (varför får hampus fler decimaler med euler än RK4. loglog plot rätt ordning).
\end{enumerate}

\section*{Uppgift 6. Titel}

\subsection*{Uppgift 6. Disa Degerholm}

\subsection*{Uppgift 6. Hampus Hillding}

\subsection*{Uppgift 6. Skillnader och likheter}

\begin{enumerate}
    \item \textbf{Analytisk lösning:} Både Hampus och Disa får att $y(t)=cost(t)$.
    \item \textbf{Hur många korrekta siffror kan man få innan datorn smälter:} Disa får ca 11 (???) korrekta siffror mha RK4 medan Hampus får 9 korrekt siffror ??? (räkna med första 3:an). Vi ser att datorn smälter då det tar för lång tid att vänta. Vi hittade sen mha mp.math, taylor och RK4 för att få 20 decimaler ???.
    \item \textbf{Hur mycket svårare blir det med skifte:} Dia härledde och insåg att k är ca 318 för att få ca t=1000. Det ger en dämpande effekt med en faktor på 300 men $t_{318}$ är betydligt svårare att få fram än $t=\pi/2$ och tar helt över med sin feltillväxt. Hampus funderar ???.
\end{enumerate}

\section*{Uppgift 7. Titel}

\subsection*{Uppgift 7. Disa Degerholm}

\subsection*{Uppgift 7. Hampus Hillding}

\subsection*{Uppgift 7. Skillnader och likheter}

\begin{enumerate}
    \item \textbf{Härled systemets ODE:} Kolla formel ???
    \item \textbf{Simulera för olika initialvärden under kort och lång tid:} Disa upptäckte att för lång tid blir det en instabil spiral utåt och för kort tid så syns inte detta. Detta är på grund av att energin ökar på grund av den numeriska härledningen ???.
    \item \textbf{Använd Semi implcit Euler för att återställa ordningen i systemet:} Check, det gick bra efterosm att ???
\end{enumerate}

\section*{Uppgift 8. Titel}

\subsection*{Uppgift 8. Disa Degerholm}

\subsection*{Uppgift 8. Hampus Hillding}

\subsection*{Uppgift 8. Skillnader och likheter}

\begin{enumerate}
    \item \textbf{Utforska numerisk för att hitta olika planetbanor s.a systemet stämmer:} Är jorden försumbar ??? 
    \item \textbf{Hur vet man att skeppet lämnar solsystemet:} Printar banan över tid och ser skillnaden på energi ??? (hur mycket)
\end{enumerate}

\end{document}